% SEGMENT OLP-0409-S01
% Part: normal-modal-logic
% Chapter: syntax-and-semantics
% Section: introduction

\documentclass[../../../include/open-logic-section]{subfiles}

\begin{document}

\olfileid{nml}{syn}{int}

\olsection{அறிமுகம்}

வாய்ப்புநிலைத் தருக்கம் \emph{வாய்ப்புநிலைக் கூற்றுகளையும்} அவற்றுக்கிடையிலான
பின்விளைவுத் தொடர்புகளையும் கையாள்கிறது. வாய்ப்புநிலைக் கூற்றுகளின் சில
எடுத்துக்காட்டுகள்:
\begin{enumerate}
\item $2+2=4$ என்பது இன்றியமையாதது.
\item நாளை மழை பெய்வது இன்றியமையாமுறையில் சாத்தியமானது.
\item $!A$ இன்றியமையாமுறையில் சாத்தியமானது எனில், $!A$ சாத்தியமானது.
\end{enumerate}
சாத்தியத்தன்மையும் இன்றியமையாமையும் மட்டுமே வாய்ப்புநிலைகள் அல்ல: பிற ஓரிட
இணைப்பிகளும் வாய்ப்புநிலைகளாக வகைப்படுத்தப்படுகின்றன. எடுத்துக்காட்டாக,
“$!A$ நிகழ வேண்டும்,” “$!A$ நிகழும்,” “$!A$ என்று டானா அறிவார்,”
அல்லது “$!A$ என்று டானா நம்புகிறார்.”

% SOURCE CORRECTION TA-NML-001: repair ordinary capitalization and missing-article defects in the historical overview.
வாய்ப்புநிலைத் தருக்கம் முதன்முதலில் அரிஸ்டாட்டிலின்
\emph{விளக்கவுரை~(De Interpretatione)} இல் தோன்றுகிறது: இன்றியமையாமை
சாத்தியத்தன்மையை உட்கிடுகிறது, ஆனால் மறுதலை அமையாது; சாத்தியத்தன்மையும்
இன்றியமையாமையும் ஒன்றையொன்று கொண்டு வரையறுக்கலாம்; $!A \land !B$
சாத்தியமாக மெய் எனில் $!A$ உம் சாத்தியமாக மெய், $!B$ உம் சாத்தியமாக
மெய், ஆனால் மறுதலை அமையாது; மேலும் $!A \to !B$ இன்றியமையாதது எனில்,
$!A$ இன்றியமையாததாக இருந்தால்~$!B$ உம் அவ்வாறே இருக்கும் என்பவற்றை
முதலில் கவனித்தவர் அவரே.

வாய்ப்புநிலைத் தருக்கத்திற்கான முதல் நவீன அணுகுமுறை சி.~ஐ. லூயிஸின்
பணியாகும்; அது லூயிஸ் மற்றும் லாங்ஃபோர்டின் \emph{குறியீட்டுத்
தருக்கம்~(Symbolic Logic)} (1932) நூலில் உச்சமடைந்தது. பொருள்சார்
நிபந்தனைவாக்கியத்தைக் கொண்டு உட்கிடையைப் பிரதிநிதித்துவப்படுத்துவதில்
லூயிஸ் \& லாங்ஃபோர்டுக்கு நிறைவு இல்லை: “$!A$ ஆனது~$!B$ ஐ உட்கிடுகிறது”
என்பதற்கு $!A \lif !B$ ஒரு குறைபாடுடைய மாற்றீடு. அதற்குப் பதிலாக,
“இன்றியமையாமுறையில், $!A$ எனில்~$!B$” என்று உட்கிடையைப்
பண்புறுத்த அவர்கள் முன்மொழிந்தனர்; இது $!A \strictif !B$ எனக்
குறிக்கப்படுகிறது. வெவ்வேறு பண்புகளை வகைப்படுத்த முயன்றபோது, லூயிஸ்
\Log{S1}, \ldots, \Log{S4}, \Log{S5} என்ற ஐந்து வெவ்வேறு வாய்ப்புநிலை
முறைமைகளை அடையாளம் கண்டார்; அவற்றின் கடைசி இரண்டு இன்னும் பயன்பாட்டில்
உள்ளன.

லூயிஸ் மற்றும் லாங்ஃபோர்டின் அணுகுமுறை முற்றிலும்
\emph{தொடரமைப்புசார்}: அவர்கள் பொருத்தமான அடிக்கோள்களையும் விதிகளையும்
அடையாளம் கண்டு, அவற்றைக் கொண்டு எவற்றை நிறுவலாம் என்று ஆராய்ந்தனர்.
ரூடால்ஃப் கார்னாப் \emph{பொருளும் இன்றியமையாமையும்~(Meaning and
Necessity)} (1947) நூலில் \emph{நிலை விவரிப்பு} என்ற கருத்தைப் பயன்படுத்தி
முதல் முயற்சி செய்யும்வரை, பொருண்மை அணுகுமுறை நீண்ட காலம் கிட்டாததாக
இருந்தது. நிலை விவரிப்பு என்பது அந்த நிலை விவரிப்பில் “மெய்யான”
அணுவாக்கியங்களின் தொகுப்பு. தன்னிச்சையான வாக்கியம்~$!A$ க்கு மெய்மை
வரையறையை விரிவாக்கியபின், எல்லா நிலை விவரிப்புகளிலும் மெய்யாக இருந்தால்
$!A$ \emph{இன்றியமையாமுறையில் மெய்} என்று கார்னாப் வரையறுக்கிறார்.
கார்னாபின் அணுகுமுறையால் \emph{அடுக்கப்பட்ட} வாய்ப்புநிலைகளைக் கையாள
முடியவில்லை; “சாத்தியமாக, இன்றியமையாமுறையில் \ldots சாத்தியமாக~$!A$”
என்ற வடிவிலுள்ள வாக்கியங்கள் எப்போதும் மிகவும் உட்புறத்திலுள்ள
வாய்ப்புநிலைக்குச் சுருங்கிவிடும்.

சால் கிரிப்கேயின் “வாய்ப்புநிலைத் தருக்கத்தில் ஒரு நிறைவுத் தேற்றம்”
(JSL 1959) என்ற கட்டுரையோடு வாய்ப்புநிலைப் பொருண்மையியலில் முக்கிய
முன்னேற்றம் ஏற்பட்டது. ஒரு கூற்று “எல்லாச் சாத்தியமான உலகங்களிலும்”
மெய்யாக இருந்தால் அது இன்றியமையாமுறையில் மெய் என்ற லைப்னிட்ஸின் எண்ணத்தை
கிரிப்கே தமது பணிக்கு அடிப்படையாகக் கொண்டார். ஆனால் ஓர் உலகம்~$w$ இல்
(அல்லது ஒரு நிலை விவரிப்பு~$s$ இல்) ஒரு கூற்றின் மெய்மை~$w$ ஐச்
சார்ந்திராததால், இந்த எண்ணமும் கார்னாபின் எண்ணத்தைப் போன்ற குறைபாட்டைக்
கொண்டுள்ளது. ஆகவே உலகங்கள் ஓர் \emph{அணுகுறவு}~$R$ ஆல் தொடர்புபட்டுள்ளன
என்று கிரிப்கே கருதினார்; “இன்றியமையாமுறையில்~$!A$” என்ற வடிவிலுள்ள
கூற்று ஓர் உலகம்~$w$ இல் மெய்யாவது, அப்போதும் அப்போது மட்டுமே,~$w$
இலிருந்து \emph{அணுகக்கூடிய} எல்லா உலகங்கள்~$w'$ இலும்~$!A$ மெய்.
இந்த அணுகுமுறையின் ஏதோ ஒரு வடிவத்தை வழங்கும் பொருண்மையியல்கள் கிரிப்கே
பொருண்மையியல் எனப்படுகின்றன; அவை வாய்ப்புநிலைத் தருக்கங்களின் பெருவளர்ச்சியை
இயலச்செய்தன.

கிரிப்கே பொருண்மையியலால் பொருள்கொள்ளப்படும்போது, \emph{தொடர்புக்
கட்டமைப்புகள்} “உள்ளிருந்து” எவ்வாறு தோன்றுகின்றன என்பதை வாய்ப்புநிலைத்
தருக்கம் நமக்குக் காட்டுகிறது. ஒரு தொடர்புக் கட்டமைப்பு என்பது ஓர் இருமத்
தொடர்பு பொருத்தப்பட்ட கணம் மட்டுமே. எடுத்துக்காட்டாக, ஒரு வகுப்பிலுள்ள
மாணவர்களின் கணத்தை அவர்களின் சமூகப் பாதுகாப்பு எண்களால் வரிசைப்படுத்தினால்,
அது ஒரு தொடர்புக் கட்டமைப்பு. ஆனால் தொடர்புக் கட்டமைப்புகள் பலவகையான
ஆட்களங்களைக் கொண்டிருக்கலாம்: உலக நிலைகளின் சார்பு சாத்தியத்தன்மையுடன்,
ஓர் ஆளின் அறிவுநிலைகள் அறிவுசார் சாத்தியத்தன்மையால் தொடர்புபட்டிருக்கலாம்;
அல்லது ஓர் இயக்க முறைமையின் நிலைகள் அவற்றின் நிலைமாற்றங்களால்
தொடர்புபட்டிருக்கலாம். இவை அனைத்தையும் வடிவமைக்க வாய்ப்புநிலைத் தருக்கத்தைப்
பயன்படுத்தலாம்: முதலாவது வழக்கமான மெய்மையியல் வாய்ப்புநிலைத் தருக்கத்தைத்
தருகிறது; மற்றவை அறிவுநிலைத் தருக்கம், இயக்கநிலைத் தருக்கம் முதலியவற்றைத்
தருகின்றன.

வாய்ப்புநிலைத் தருக்கவியலாளர்களுக்கு “ஒத்திசைவுக் கோட்பாடு” என்று
அறியப்படும் ஒரு குறிப்பிட்ட கோணத்தில் நாம் கவனம் செலுத்துகிறோம்.
அணுகுறவு~$R$ இன் பல பண்புகளை (அது கடப்புநிலை உடையதா, சமச்சீரானதா
முதலியவற்றை) பொருத்தமான “வாய்ப்புநிலைத் திட்டவடிவங்களைக்” கொண்டு
\emph{வாய்ப்புநிலை மொழியிலேயே} பண்புறுத்தலாம் என்பது கிரிப்கேயின் மிகவும்
முக்கியமான தொடக்கக் கண்டுபிடிப்புகளில் ஒன்று. எடுத்துக்காட்டாக, $R$ இன்
தற்பின்னோக்கிய தன்மை “இன்றியமையாமுறையில்~$!A$ எனில்,~$!A$” என்ற
திட்டவடிவத்திற்கு “ஒத்துள்ளது” என்று வாய்ப்புநிலைத் தருக்கவியலாளர்கள்
கூறுகின்றனர். \Ax{D}, \Ax{T}, \Ax{B}, \Ax{4}, \Ax{5} ஆகிய
திட்டவடிவங்களின் சேர்க்கையால் பெறப்படும் பல செவ்வியல் வாய்ப்புநிலைத் தருக்க
முறைமைகளின் (எடுத்துக்காட்டாக, \Log{S4}, \Log{S5}) ஒத்திசைவுக்
கோட்பாட்டையே நாம் முக்கியமாக ஆராய்கிறோம்.

\end{document}
